\section{Preliminaries}
\label{sec:prelim}

Let $q=p^s$ for some prime $p$ and $s\in \bbN$. Let $\f:\bbF^n\rightarrow\bbF_q$ be a function over the finite field $\bbF_q$. Let $\omega_p=e^{2\pi i/p}$ be the root of unity modolu $p$. 

\begin{definition}[Trace function in finite fields]
The trace function $\Tr:\bbF_q\rightarrow\bbF_p$ is defined as follows:
$$\Tr(x) = \sum_{i=0}^{s-1}x^{p^i}$$
\end{definition}

% \begin{remark}
% Observe that for $s=1$, meaning that $q=p$, the trace function is simply the identity function. 
% \end{remark}

\begin{definition}
The character group of $\bbF_q$, denoted \characters, contains all homomorphisms $\chi:\bbF^n\rightarrow \bbC$.
\end{definition}

\begin{fact}
\characters consists of all the functions:
$ \chi_{\alpha}(x) = \omega_p^{\Tr(\langle \alpha, x\rangle)}
 $, for all $\alpha\in\bbF^n$. 
\end{fact}

\begin{fact}
$\{\chi_\alpha\}_{\alpha\in \bbF^n}$ form an orthonormal basis for all functions $\f : \bbF^n_q \rightarrow \bbC$ under the inner-product $\langle \f,\g\rangle = \bbE_{x\in \bbF^n}[f(x)\overline{g(x)}].$
\end{fact}


\begin{definition}
Every function  $h:\bbF^n\rightarrow\bbC$ has a Fourier expansion $\hat{h}:\bbF^n\rightarrow\bbC$ given by
$$
h(x)=\sum_{\alpha\in\bbF_q}\hat{h}(\alpha)\chi_{\alpha}(x)
$$
$$
\hat{h}(\alpha)=q^{-n}\sum_{x\in\bbF_q}{h}(x)\overline{\chi_{\alpha}(x)}
$$
\end{definition}

\begin{definition}
The norm of a function  $h:\bbF^n\rightarrow\bbC$ is equal to:
\[
\norm{h}^2=\inner{h}{h}=\bbE_{x\in \bbF^n}[h(x)\overline{h(x)}]=q^{-n}\sum_{x\in\bbF^n}h(x)^2
\]
\end{definition}

\begin{corollary}
Let $h:\bbF^n\rightarrow\bbC$. Then: $\norm{h}\le1$. 
\end{corollary}
\begin{proof}
$$
\norm{h}^2= q^{-n}\sum_{x\in\bbF^n}h(x)^2\le q^{-n}\sum_{x\in\bbF^n}|h(x)|^2\le 1
$$
\end{proof}

\begin{theorem}Let $h_1,h_2:\bbF^n\rightarrow\bbC$.
\begin{itemize}
    \item Parseval's Identity: $\norm{h_1}^2 = \sum_{\alpha\in\bbF^n}\widehat{h_1}(\alpha)^2$.
    \item Plancheral's Identity: $\inner{h_1}{h_2} = \sum_{\alpha\in\bbF^n}\widehat{h_1}(\alpha)\widehat{h_2}(\alpha)$.
\end{itemize}
\end{theorem}

As there is no isomorphism between $\bbF_q$ and $\bbC$, we instead consider $q-1$ functions $\ff_c:\bbF^n\rightarrow \bbC$ for all $c\in \bbF_q^{*}$:
\begin{align}
\ff_c(x) := \omega_p^{\Tr(c\cdot f(x))} = \sum_{\alpha\in \bbF^n} \widehat{\left(\ff_c\right)}(\alpha)\chi_{\alpha}(x)
\end{align}
\begin{definition}
For any $\f:\bbF^n\rightarrow \bbF_q$ we defined the $\bbF_p$-expansion of $f$ as the set:
$
\left\{\ff_c\right\}_{c\in\bbF_q^{*}}
$.
\end{definition}

\begin{definition}
Let $\f,\g:\bbF^n\rightarrow\bbF_q$. We defined the distance between $f,g$ as:
\[
\delta(\f,\g) =\Pr_{x\in\bbF^n}\left[f(x)\ne g(x)\right] 
\]
\end{definition}

Now, we express the distance of $\f,\g$ in terms of the Fourier coefficients $\bbF_p$-expansion sets. 
\begin{claim} \label{claim: delta(f,g)}Let $f,g:\bbF^n\rightarrow\bbF_q$ and their $\bbF_p$-expansion sets be $\{\fexp{\ff}{c}\}_{c\in\bbF_q^{*}},\{\fexp{\fg}{c}\}_{c\in\bbF_q^{*}}$ respectively.
$$\delta(f,g) = 1-\frac{1}{q}\left(1+\sum_{c\in\bbF_q^{*}}\inner{\fexp{\ff}{c}}{\fexp{\fg}{c}}\right) = 1-\frac{1}{q}\left(1+\sum_{c\in\bbF_q^{*}}\sum_{\alpha\in\bbF^n}\widehat{\fexp{\ff}{c}}(\alpha)\overline{\widehat{\fexp{\fg}{c}}(\alpha)} \right)$$    
\end{claim}
\begin{definition}
Let $\lin$ be the set of linear functions form $\bbF^n$ to $\bbF_q$. For $\alpha\in\bbF^n$, we denote $\ell_{\alpha}(x) = \inner{x}{\alpha}$. Then:
\[
\lin = \left\{\ell_{\alpha}\right\}_{\alpha\in\bbF^n}
\]
\end{definition}

\begin{claim}
Let $\f:\bbF^n\rightarrow\bbF_q$ with $\bbF_p$-expansion sets be $\{\fexp{\ff}{c}\}_{c\in\bbF_q^{*}}$ and let $\ell_{\alpha}\in\lin$ for some $\alpha\in\bbF^n$, and let it's $\bbF_p$-expansion set be $\{\fexp{\lambda}{c}\}_{c\in\bbF_q^{*}}$. Then:
$$
\delta(\f,\ell_{\alpha}) = 1-\left(\frac{1}{q}\left(1+\sum_{c\in\bbF_q^{*}}\widehat{\fexp{\ff}{c}}(c\alpha)\right)\right)  
$$
\end{claim}
\begin{proof} For all $\beta\in\bbF^n$, the following is 0 if $\beta \ne c\alpha$ and 1 otherwise:
$$\widehat{\fexp{\lambda}{c}}(\beta)=q^{-n}\sum_{x\in\bbF^n}\fexp{\lambda}{c}(x)\overline{\chi_{\beta}(x) }= q^{-n}\sum_{x\in\bbF^n}\omega_p^{\Tr(c\cdot \ell_\alpha(x)) - \Tr(\langle \beta, x\rangle) }=q^{-n}\sum_{x\in\bbF^n}\omega_p^{\Tr(\langle c\alpha, x\rangle) - \Tr(\langle \beta, x\rangle) }$$ 
\end{proof}

\begin{corollary}\label{bound for sum of expansion q}For all $\alpha\in \bbF^n$, 
$-1\le \sum_{c\in\bbF_q^{*}}\widehat{\fexp{\ff}{c}}(c\alpha)\le q-1$ is a real number.
\end{corollary}
\begin{proof}
This follows because $\delta(f,\ell_{\alpha})\in[0,1]$.
\end{proof}

\begin{proof}
    Equivalent to $0\le\delta\le 1-1/q$
\end{proof}

\begin{definition}
The distance of $\f:\bbF^n\rightarrow\bbF_q$ from \lin is defined as:  $\delta(f,\lin) :=\min_{\alpha\in\bbF^n}\delta(f,\ell_\alpha)$.
\end{definition}

\begin{corollary}\label{distance from linear}
\[
\delta(f,\lin) = \min_{\alpha\in\bbF^n}\left(1-\left(\frac{1}{q}\left(1+\sum_{c\in\bbF_q^{*}}\widehat{\fexp{\ff}{c}}(c\alpha)\right)\right) \right)
\]
Equivalently,
$$\delta(f, \mathcal{LIN}) = 1- \frac{1}{q}\left(1+ \max_{\alpha\in\bbF^n}\left(\sum_{c\in\bbF_q^{*}}\widehat{\fexp{\ff}{c}}(c\alpha)\right)\right)$$
\end{corollary}



%%%%%%%%%%%%%%%%%%%%%%%%%%%%%%%%%%%%%%%%%%%%%%%%%%%%%%%%%%%%%%%%%%%%%%%%%%%%%%%%
%%%%%%%%%%%%%%%%%%%%%%%%%%%%%%%%%%%%%%%%%%%%%%%%%%%%%%%%%%%%%%%%%%%%%%%%%%%%%%%%
%%%%%%%%%%%%%%%%%%%%%%%%%%%%%%%%%%%%%%%%%%%%%%%%%%%%%%%%%%%%%%%%%%%%%%%%%%%%%%%%
\section{BLR Test over Finite Fields}
\label{sec:blr over finite fields}

\paragraph{The BLR Test}Given oracle access to a function $\f:\bbF^n\rightarrow\bbF_q$, the verifier performs the following:
\begin{itemize}
    \item Sample random $a,b\in\bbF^n\setminus\{0\}$ and sample $x,y\in\bbF^n$. 
    \item $\BLRVerifier=1 \ \iff \ a\cdot f(x)+b\cdot f(y) = f(a\cdot x+b\cdot y)$. 
\end{itemize}

\begin{theorem}\label{FieldBLR}
Let $\f:\bbF^n\rightarrow\bbF_q$ with $\bbF_p$-expansion sets be $\{\fexp{\ff}{c}\}_{c\in\bbF_q^{*}}$. 

\[
\Pr_{a,b,x,y}\left[\BLRVerifier=1 \right] = \frac{1}{q}\left(1+\frac{1}{(q-1)^2}\sum_{\alpha\in\bbF^n}\left(\sum_{a\in\bbF_q^{*}}\widehat{\fexp{\ff}{a}}(a\alpha)\right)^3\right)
\]
\end{theorem}


\begin{proof}
Let the $\bbF_p$-expansion set of $f$ be $\{\fexp{\ff}{c}\}_{c\in\bbF_q^{*}}$.
Let also $a,b\in\bbF^n\setminus\{0\}$ and $x,y\in\bbF^n$. We observe that:
\[
\BLRVerifier=1  \ \iff \ \frac{1}{q}\left(\sum_{c\in\bbF_q}\omega_p^{\Tr(c(af(x)+bf(y)-f(ax+by)}\right) = 1
\]
As the above sum is zero unless $af(x)+bf(y)-f(ax+by) = 0$. Hence:
\begin{align*}
\Pr_{a,b,x,y}\left[\BLRVerifier=1\right]&=\bbE_{a,b,x,y}\left[\frac{1}{q}\left(\sum_{c\in\bbF_q}\omega_p^{\Tr(c(af(x)+bf(y)-f(ax+by))}\right)\right]\\
&= \frac{1}{q}+\bbE_{a,b,x,y}\left[\sum_{c\in\bbF_q^{*}}\omega_p^{\Tr(c(af(x)+bf(y)-f(ax+by))}\right]
\end{align*}

Now, we consider the last term:
\begin{align*}
&\bbE_{a,b,x,y}\left[\sum_{c\in\bbF_q^{*}}\omega_p^{\Tr(c(af(x)+bf(y)-f(ax+by))}\right] \\&= \bbE_{a,b,x,y}\left[\sum_{c\in\bbF_q^{*}}\fexp{\ff}{ca}(x)\fexp{\ff}{cb}(y)\fexp{\ff}{-c}(a\cdot x+b\cdot y)\right]\\
&=\bbE_{a,b,x,y}\left[\sum_{c\in\bbF_q^{*}}\sum_{\alpha,\beta,\gamma\in\bbF^n}\widehat{\fexp{\ff}{ca}}(\alpha)\chi_{\alpha}(x)\widehat{\fexp{\ff}{cb}}(\beta)\chi_{\beta}(y)\widehat{\fexp{\ff}{-c}}(\gamma)\chi_{\gamma}(ax+by)\right]\\
&= \bbE_{a,b,x,y}\left[\sum_{c\in\bbF_q^{*}}\sum_{\alpha,\beta,\gamma\in\bbF^n}\widehat{\fexp{\ff}{ca}}(\alpha)\widehat{\fexp{\ff}{cb}}(\beta)\widehat{\fexp{\ff}{-c}}(\gamma)\chi_{\alpha}(x)\chi_{\beta}(y)\chi_{\gamma}(ax+by)\right]\\
&= \bbE_{a,b,x,y}\left[\sum_{c\in\bbF_q^{*}}\sum_{\alpha,\beta,\gamma\in\bbF^n}\widehat{\fexp{\ff}{ca}}(\alpha)\widehat{\fexp{\ff}{cb}}(\beta)\widehat{\fexp{\ff}{-c}}(\gamma)\chi_{\alpha}(x)\chi_{\beta}(y)\chi_{a\gamma}(x)\chi_{b\gamma}(y)\right]\\
&= \bbE_{a,b}\left[\sum_{c\in\bbF_q^{*}}\sum_{\alpha,\beta,\gamma\in\bbF^n}\widehat{\fexp{\ff}{ca}}(\alpha)\widehat{\fexp{\ff}{cb}}(\beta)\widehat{\fexp{\ff}{-c}}(\gamma)\bbE_{x,y}\left[\chi_{\alpha}(x)\chi_{\beta}(y)\chi_{a\gamma}(x)\chi_{b\gamma}(y)\right]\right]
\end{align*}

Now, we consider the following expression:
\begin{align*}
    \bbE_{x,y}\left[\chi_{\alpha}(x)\chi_{\beta}(y)\chi_{a\gamma}(x)\chi_{b\gamma}(y)\right] = \bbE_{x,y}\left[\omega_p^{\Tr(\inner{\alpha}{x}+\inner{\beta}{y}+ \inner{\gamma}{ax}+\inner{\gamma}{by})}\right]
\end{align*}
Observe that the above is zero, unless $\gamma = -a^{-1}\alpha$ and $\beta = a^{-1}b\alpha$. Thus:
\begingroup
\allowdisplaybreaks
\begin{align*}
&\bbE_{a,b,x,y}\left[\sum_{c\in\bbF_q^{*}}\omega_p^{\Tr(c(af(x)+bf(y)-f(ax+by))}\right]\\
&=
\bbE_{a,b}\left[\sum_{c\in\bbF_q^{*}}\sum_{\alpha,\beta,\gamma\in\bbF^n}\widehat{\fexp{\ff}{ca}}(\alpha)\widehat{\fexp{\ff}{cb}}(\beta)\widehat{\fexp{\ff}{-c}}(\gamma)\bbE_{x,y}\left[\chi_{\alpha}(x)\chi_{\beta}(y)\chi_{a\gamma}(x)\chi_{b\gamma}(y)\right]\right] \\
&= 
\bbE_{a,b}\left[\sum_{c\in\bbF_q^{*}}\sum_{\alpha\in\bbF^n}\widehat{\fexp{\ff}{ca}}(\alpha)\widehat{\fexp{\ff}{cb}}(a^{-1}b\alpha)\widehat{\fexp{\ff}{-c}}(-a^{-1}\alpha)\right]\\
&=
\bbE_{a,b}\left[\sum_{c\in\bbF_q^{*}}\sum_{\alpha\in\bbF^n}\widehat{\fexp{\ff}{ca}}(ca\alpha)\widehat{\fexp{\ff}{cb}}(cb\alpha)\widehat{\fexp{\ff}{-c}}(-c\alpha)\right]\\
&= \bbE_{a,b}\left[\sum_{c\in\bbF_q^{*}}\sum_{\alpha\in\bbF^n}\widehat{\fexp{\ff}{a}}(a\alpha)\widehat{\fexp{\ff}{b}}(b\alpha)\widehat{\fexp{\ff}{c}}(c\alpha)\right]\\
&= \frac{1}{(q-1)^2}\sum_{a,b,c\in\bbF_q^{*}}\sum_{\alpha\in\bbF^n}\widehat{\fexp{\ff}{a}}(a\alpha)\widehat{\fexp{\ff}{b}}(b\alpha)\widehat{\fexp{\ff}{c}}(c\alpha)\\
&= \frac{1}{(q-1)^2}\sum_{\alpha\in\bbF^n}\left(\sum_{a\in\bbF_q^{*}}\widehat{\fexp{\ff}{a}}(a\alpha)\right)^3
\end{align*}
\endgroup

\end{proof}


\begin{theorem}
Let $\f:\bbF^n\rightarrow\bbF_q$. 
\[
(1-\delta(f,\lin))^3 \le \Pr_{a,b,x,y}\left[\BLRVerifier=1\right] \le 1-\delta(f,\lin)
\]
\end{theorem}
\begin{proof}
For the left hand-side:
\begin{align*}
\Pr\left[\BLRVerifier=1 \right] &= \frac{1}{q}\left(1+\frac{1}{(q-1)^2}\sum_{\alpha\in\bbF^n}\left(\sum_{a\in\bbF_q^{*}}\widehat{\fexp{\ff}{a}}(a\alpha)\right)^3\right)\\
&\le  \frac{1}{q}\left(1+\frac{1}{(q-1)^2}\max_{\alpha\in\bbF^n}\left(\sum_{a\in\bbF_q^{*}}\widehat{\fexp{\ff}{a}}(a\alpha)\right)\sum_{\alpha\in\bbF^n}\left(\sum_{a\in\bbF_q^{*}}\widehat{\fexp{\ff}{a}}(a\alpha)\right)^2\right)\\
&=  \frac{1}{q}\left(1+\frac{1}{(q-1)^2}\max_{\alpha\in\bbF^n}\left(\sum_{a\in\bbF_q^{*}}\widehat{\fexp{\ff}{a}}(a\alpha)\right)\sum_{\alpha\in\bbF^n}\left(\sum_{a,b\in\bbF_q^{*}}\widehat{\fexp{\ff}{a}}(a\alpha)\widehat{\fexp{\ff}{b}}(b\alpha)\right)\right)\\
&=  \frac{1}{q}\left(1+\frac{1}{(q-1)^2}\max_{\alpha\in\bbF^n}\left(\sum_{a\in\bbF_q^{*}}\widehat{\fexp{\ff}{a}}(a\alpha)\right)\sum_{a,b\in\bbF_q^{*}}\left(\sum_{\alpha\in\bbF^n}\widehat{\fexp{\ff}{a}}(a\alpha)\widehat{\fexp{\ff}{b}}(b\alpha)\right)\right)
\end{align*}
Consider the following, due to Plancherel's identity and the fact that $\fexp{\ff}{a}$ map to the roots of unity modulo $p$:
\begin{align*}
\sum_{\alpha\in\bbF^n}\widehat{\fexp{\ff}{a}}(a\alpha)\widehat{\fexp{\ff}{b}}(b\alpha) 
= \bbE_{\alpha}\left[{\fexp{\ff}{a}(a\alpha)\fexp{\fg}{b}(b\alpha)}\right] 
\le \bbE_{\alpha}\left[{\abs{\fexp{\ff}{a}(a\alpha)\fexp{\fg}{b}(b\alpha)}}\right]\le 1
\end{align*}

Hence, it follows that:
\[
 \Pr_{a,b,x,y}\left[\BLRVerifier=1\right] \le   \frac{1}{q}\left(1+\max_{\alpha\in\bbF^n}\left(\sum_{a\in\bbF_q^{*}}\widehat{\fexp{\ff}{a}}(a\alpha)\right)\right)
= 1-\delta(f,\lin)
\]

For the left hand side of the claim:
\begin{align*}
\Pr\left[\BLRVerifier=1 \right] &= \frac{1}{q}\left(1+\frac{1}{(q-1)^2}\sum_{\alpha\in\bbF^n}\left(\sum_{a\in\bbF_q^{*}}\widehat{\fexp{\ff}{a}}(a\alpha)\right)^3\right)\\
&\ge  \frac{1}{q}\left(1+\frac{1}{(q-1)^2}\max_{\alpha\in\bbF^n}\left(\sum_{a\in\bbF_q^{*}}\widehat{\fexp{\ff}{a}}(a\alpha)\right)^3\right)\\
&\ge  \left(\frac{1}{q}\left(1+\max_{\alpha\in\bbF^n}\left(\sum_{a\in\bbF_q^{*}}\widehat{\fexp{\ff}{a}}(a\alpha)\right)\right)\right)^3\\
&=  \left(1-\delta(f,\lin)\right)^3
\end{align*}
where the last inequality follows from the following: for every $q\geq 2, A\ge \frac{1-|\bbF|^2}{2q-1}$:
\[
\left(\frac{1}{q}\left(1+A\right)\right)^3\le
\frac{1}{q}\left(1+\frac{1}{(q-1)^2}A^3\right)
\]
(we plug $A=\max_{\alpha\in\bbF^n}\left(\sum_{a\in\bbF_q^{*}}\widehat{\fexp{\ff}{a}}(a\alpha)\right)$ in the expression above. By \cref{bound for sum of expansion q}, we observe that the above holds for any $q$).

\end{proof}




% \inputleanmodule{BlrPcp.Basic}

\section{previous chatgpt generated content}
\begin{definition}
    \label{def:ff-linear}
    A function $f : F^n \to F$ is linear if there exists
    $a \in F^n$ such that
    \[
        f(x) = \sum_i a_i x_i
    \]
    for every $x \in F^n$.
\end{definition}

\begin{definition}
    \label{def:ff-blr-test}
    The finite-field BLR test samples $x,y \in F^n$ and
  $a,b \in F \setminus \{0\}$, and accepts iff
  \[
    a f(x) + b f(y) = f(ax + by).
  \]
\end{definition}

\begin{definition}
    \label{def:ff-blr-reject-prob}
    The rejection probability of the finite-field BLR test is
  \[
    \Pr[V_{BLR}^f = 0] = 1 - \Pr[V_{BLR}^f = 1].
  \]
\end{definition}

\begin{definition}
    \label{def:ff-additive-characters}
    Write $F=\mathbb{F}_q$ where $q = p^m$ for a prime number $p$ and a positive integer $m$. For each $\alpha \in \mathbb{F}_q^n$, define
  the additive character $\chi_\alpha : \mathbb{F}_q \rightarrow \mathbb{C}$ as
  \[
    \chi_\alpha(x) := \omega_p^{\mathrm{Tr}(\langle \alpha, x\rangle)},
  \]
  where $ \omega_p = \exp(2π i/p)$ and  $\mathrm{Tr} : \mathbb{F}_q \to \mathbb{F}_p $ is the field trace $\mathrm{Tr}(a) = \sum_{j=0}^{m-1} a^{p^j}$.
\end{definition}

\begin{lemma}
    \label{lem:ff-characters-orthonormal}
The family $\{\chi_\alpha\}_{\alpha \in \mathbb{F}_q^n}$ is an
  orthonormal basis for the complex-valued functions
  $g : \mathbb{F}_q^n \to \mathbb{C}$ with respect to the inner product
  \[
    \langle g,h\rangle := \mathbb{E}_{x \in \mathbb{F}_q^n}[g(x)\overline{h(x)}].
  \]
\end{lemma}

\begin{lemma}
    \label{lem:ff-fourier-expansion}
Every function $g : \mathbb{F}_q^n \to \mathbb{C}$ has a unique
  Fourier expansion
  \[
    g(x) = \sum_{\alpha \in \mathbb{F}_q^n} \widehat g(\alpha)\chi_\alpha(x).
  \]
\end{lemma}

\begin{lemma}
    \label{lem:ff-parseval-plancherel}
    Parseval and Plancherel hold:
  \[
    \langle g,g\rangle = \sum_\alpha |\widehat g(\alpha)|^2,
    \qquad
    \langle g,h\rangle = \sum_\alpha \widehat g(\alpha)\overline{\widehat h(\alpha)}.
  \]
\end{lemma}

\begin{definition}
    \label{def:ff-fourier-set}
    For $f : \mathbb{F}_q^n \to \mathbb{F}_q$, define its Fourier set
  \[
    \Phi(f) := \{\varphi_c\}_{c \in \mathbb{F}_q^\times},
    \qquad
    \varphi_c(x) := \omega_p^{\mathrm{Tr}(c f(x))}.
  \]
\end{definition}

\begin{lemma}
    \label{lem:ff-distance-formula}
If $f,g : \mathbb{F}_q^n \to \mathbb{F}_q$ have Fourier sets
  $\{\varphi_c\}_{c \in \mathbb{F}_q^\times}$ and
  $\{\psi_c\}_{c \in \mathbb{F}_q^\times}$, then
  \[
    \Delta(f,g)
    = 1 - \frac{1}{q}\left(1 + \sum_{c \in \mathbb{F}_q^\times}
      \langle \varphi_c,\psi_c\rangle\right).
  \]
\end{lemma}

\begin{lemma}
    \label{lem:ff-distance-to-linearity}
    If $f$ has Fourier set $\{\varphi_c\}_{c \in \mathbb{F}_q^\times}$,
  then its distance to the linear functions is
  \[
    \Delta(f,\mathrm{LIN})
    = 1 - \frac{1}{q}\left(1 + \max_{\alpha \in \mathbb{F}_q^n}
      \sum_{c \in \mathbb{F}_q^\times} \widehat{\varphi_c}(c\alpha)\right).
  \]
\end{lemma}

\begin{lemma}
    \label{lem:ff-spectral-bound}
A Fourier-analytic calculation shows that the rejection probability
  of the finite-field BLR test dominates the Fourier-set expression for
  $\Delta(f,\mathrm{LIN})$
\end{lemma}

\begin{lemma}
    \label{thm:ff-blr-improved-analysis}
For every function $f : F^n \to F$, the finite-field BLR test
  satisfies
  \[
    \Pr[V_{BLR}^f = 0] \ge \Delta(f,\mathrm{LIN}).
  \]
  This is the improved finite-field analysis proved via Fourier analysis.
\end{lemma}

\begin{lemma}
    \label{cor:ff-high-acceptance-close-to-linear}
If the finite-field BLR test accepts $f$ with probability at least
  $1 - \varepsilon$, then
  \[
    \Delta(f,\mathrm{LIN}) \le \varepsilon.
  \]
\end{lemma}

\begin{lemma}
    \label{rem:q2-special-case}
    When $q = 2$, this reduces to the Boolean Fourier-analytic proof
  of the BLR theorem.
\end{lemma}